\chapter{General one-body hopping matrices}
\label{chap:generalhopping}

The nearest-neighbor chain is a specialization of
\begin{equation}
H(h)=\sum_{ij,\sigma}h_{ij}c_{i\sigma}^\dagger c_{j\sigma}
 +U\sum_i n_{i\uparrow}n_{i\downarrow}.
\label{eq:generalh}
\end{equation}
The interface convention is explicit: \code{h[i,j]} is the coefficient of
$c_i^\dagger c_j$.  The same $h$ is applied to both spin species.

\section{Validation and supported terms}

The input must have shape $(L,L)$, contain finite numerical values, and satisfy
\begin{equation}
h=h^\dagger
\end{equation}
to relative and absolute tolerance $10^{-12}$.  Accepted roundoff-level
asymmetry is projected away by storing $(h+h^\dagger)/2$.  The validated matrix
is copied and made read-only, preventing caller mutation from changing a
calculation after construction.

The entries can describe
\begin{itemize}
\item arbitrary-range or graph hopping;
\item real or complex Peierls phases;
\item spin-independent onsite potentials on the diagonal;
\item dense matrices produced by a one-particle basis transformation.
\end{itemize}
Each nonzero entry is evaluated through the same $c_i^\dagger c_j$ operator
path used earlier.  CSR and matrix-free complex-Hermitian actions are therefore
available without introducing a second sign convention.

For a ring threaded by total phase $\Phi$, one convenient gauge is
\begin{equation}
h_{i+1,i}=-t e^{i\Phi/L},\qquad
h_{i,i+1}=\left(h_{i+1,i}\right)^*.
\end{equation}

\lstinputlisting[caption={A complex flux model with longer-range hopping.},label={lst:generalh}]{code/step07_generic_hopping.py}

Complex Hermitian matrices produce real eigenvalues but generally complex
eigenvectors.  The generalized solver therefore preserves complex vector
dtype, uses a deterministic complex initial vector, and recomputes complex
residuals with the matrix-free action.

\section{The orbital-rotation caveat}

The general one-body interface is necessary for wavelets and orbital rotations,
but it is not sufficient for the interacting model.  Let a unitary matrix $R$
define new operators through
\begin{equation}
c_{i\sigma}=\sum_a R_{ia}d_{a\sigma}.
\end{equation}
The one-body matrix transforms simply:
\begin{equation}
h' = R^\dagger h R.
\end{equation}
The onsite interaction does not generally remain onsite.  Instead,
\begin{align}
U\sum_i n_{i\uparrow}n_{i\downarrow}
&=\sum_{abcd}V_{abcd}
d_{a\uparrow}^\dagger d_{b\uparrow}
d_{c\downarrow}^\dagger d_{d\downarrow},\\
V_{abcd}
&=U\sum_i R^*_{ia}R_{ib}R^*_{ic}R_{id}.
\label{eq:rotatedinteraction}
\end{align}
Therefore, passing only $h'$ to Eq.~\eqref{eq:generalh} while retaining a local
$U$ usually defines a different physical model.  It is exact at $U=0$ and in
special cases where the rotation preserves the local interaction.  A general
four-index two-body operator is the required next layer for faithful wavelet,
natural-orbital, or learned-orbital calculations at finite $U$.

\paragraph{Checkpoint.}
At $U=0$, compare the complete many-body spectrum obtained from a general $h$
with sums of its one-particle eigenvalues.  For complex $h$, additionally test
Hermiticity, CSR versus matrix-free action on complex random vectors, real
eigenvalues, and small residuals.
